%% LyX 2.0.7.1 created this file.  For more info, see http://www.lyx.org/.
%% Do not edit unless you really know what you are doing.
\documentclass[11pt,english]{article}
\usepackage[latin9]{inputenc}
\usepackage{geometry}
\geometry{verbose,tmargin=1in,bmargin=1in,lmargin=1in,rmargin=1in}
\usepackage{varioref}
\usepackage{float}
\usepackage{textcomp}
\usepackage{amstext}
\usepackage{amssymb}
\usepackage{graphicx}

\makeatletter

%%%%%%%%%%%%%%%%%%%%%%%%%%%%%% LyX specific LaTeX commands.
\newcommand{\lyxmathsym}[1]{\ifmmode\begingroup\def\b@ld{bold}
  \text{\ifx\math@version\b@ld\bfseries\fi#1}\endgroup\else#1\fi}

%% A simple dot to overcome graphicx limitations
\newcommand{\lyxdot}{.}


\makeatother

\usepackage{babel}
\begin{document}

\title{Mapping The Warp of the Milky Way Galaxy}


\author{Patrick Kantorski}


\date{May 5, 2014}
\maketitle
\begin{abstract}
Utilizing 3 hours of observation time, my lab parters and I collected
data on the power spectrum of the High Velocity Clouds in the Milky
Way galaxy using the radio telescope at the Leuschner Observatory.
Our pointing ranged from $l=75^{\lyxmathsym{\textopenbullet}}$, $b=+20^{\lyxmathsym{\textopenbullet}}\rightarrow-20^{\text{\textopenbullet}}$
with each pointing collected with an integration time of about 4 minutes.
After converting frequency information into velocities and distances,
we convolved a Gaussian profile as the kernel to smooth the data and
also removed any noise below temperatures of 0K. Using the Python
module 'PyLab', we created an image plotting the distance from and
height above the galactic center, mapping temperatures to color intensities.
The result of the experiment shows an uneven distribution of H1 emissions
above and below the galactic plane.
\end{abstract}
\tableofcontents{}


\section{Introduction}


\subsection{The Leuschner Radio Telescope}

The radio telescope used for in this lab is located at the Leuschner
Observatory at a latitude of $+37.8732^{\lyxmathsym{\textopenbullet}}$
and longitude of $\lyxmathsym{\textminus}122^{\lyxmathsym{\textopenbullet}}9\lyxmathsym{\textasciiacute}$
in Berkeley, California. The Leuschner Observatory is a student observatory
ran by the Astronomy Department of the University of California at
Berkeley. The diameter of the dish is approximately 4.5m and operates
at a frequency range from 1320\textminus{}1740MHz with a 110MHz bandwidth.
The telescope operates with mixed local oscillator signals transferred
to an ADC after passing through two bandpass filters in order to minimize
contaminations from other signal frequencies.

\begin{figure}[H]
\begin{centering}
\includegraphics[scale=0.48]{/Users/ppkantorski/Documents/Radio_Astronomy/Lab_4/Pictures/Leuschner_circuit}
\par\end{centering}

\caption{This is a picture of the circuit diagram schematics for the 4.5m diameter
radio telescope at the Leuschner Observatory in Berkeley, California.
(Image was taken from http://ugastro.berkeley.edu/radio/leuschner/Leuschner\_circuit.png
on May 4, 2014.) \label{fig: Doppler Shift-1-2}}
\end{figure}
\begin{figure}[H]
\begin{centering}
\includegraphics[scale=0.18]{/Users/ppkantorski/Documents/Radio_Astronomy/Lab_4/Pictures/dish}
\par\end{centering}

\caption{This is a picture of the 4.5m diameter radio telescope at the Leuschner
Observatory in Berkeley, California ($\lyxmathsym{\textminus}122^{\lyxmathsym{\textopenbullet}}9^{\prime}$,$+37.8732^{\lyxmathsym{\textopenbullet}}$).
Wiring runs from the telescope to the near by building underground
where they are processed through the electronics required for operation.
(Image was taken from http://ugastro.berkeley.edu/pictures/science/leuschner/radio/dish.jpg
on May 4, 2014.) \label{fig: Doppler Shift-1}}
\end{figure}
\begin{figure}[H]
\begin{centering}
\includegraphics[scale=0.44]{/Users/ppkantorski/Documents/Radio_Astronomy/Lab_4/Pictures/feedmatch_equipment_writing}
\par\end{centering}

\caption{This is a picture of the electronics used to operate the 4.5m radio
telescope at the Leuschner Observatory along with the actual reflector
detector of the telescope itself. (Image was taken from http://ugastro.berkeley.edu/pictures/science/leuschner/radio/feedmatch\_equipment\_writing.jpg
on May 4, 2014.) \label{fig: Doppler Shift-1-1}}
\end{figure}



\subsection{Mapping The Galactic Warp}

External galaxies are often distorted by nearby galactic neighbors.
To determine if the Milky Way galaxy is distorted or 'warped', my
partners and I used the radio telescope at the Leuschner Observatory
and profiled a section of the sky vertically slicing across the galactic
plane. Hypothetically, a 'flat' galaxy would have emission peaks everywhere
at galactic latitude $b=0^{\circ}$. To observe and display these
distortions, we were only able to take one complete vertical profile
of the Milky Way. If our galaxy does experience a warp, the result
of this experiment should show an uneven distribution of H1-emissions
along galactic latitudes above and below $b=0^{\circ}$.


\section{Background Theory}


\subsection{Hydrogen 21cm Line}

The intrinsic spin property of the electron in a Hydrogen causes the
atom to emit a photon at precisely a distance known as the 21cm line
or a frequency of 1420.4MHz. Additionally, the transition between
the three-fold degenerate triplet states and the non-degenerate singlet
state in Hydrogen atoms have an energy difference $\Delta E=h\nu\thickapprox5.874\times10^{-9}eV$.
This energy difference in transmitted by photons when the spin state
of the electron flip from their alignment with the spin state of the
proton to an anti-aligned position. Since the frequency of 1420.4MHz
is in the radio wave band of the light spectrum, H1-emissions penetrate
most of the dust along the line-of-sight and are not absorbed or scattered.
Utilizing the 21cm line, radio astronomers can examine spectra line
emissions and study the propagation of H1-emitted photons across the
vacuum of space.


\subsection{The Doppler Effect}

The Doppler effect is the change in frequency of a wave or periodic
event for an observer in motion relative to a source. An everyday
example of the doppler effect can be examined when a vehicle emitting
a siren or horn approaches, passes, and then recedes away from an
observer. As the sound waves approach, the frequency of the sound
increases, making the pitch of the sound also higher. As the vehicle
recedes, the frequency and pitch diminishes. This is because each
wave takes slightly less time to reach the observer than the previous.
For light waves, the effect is the same. While photons travel through
space, the color or frequency of the wave also stretches and contracts.
making the waves \char`\"{}bunch together\char`\"{}. On the contrary,
if the wave source is moving away from the observer, the waves \char`\"{}spread
out\char`\"{}. This ``bunching'' and ``spreading'' of propagating
light waves is known as red and blue shift.

\begin{figure}[H]
\begin{centering}

\includegraphics[scale=0.46]{/Users/ppkantorski/Documents/Radio_Astronomy/Lab_4/Pictures/doppler_effect}
\par\end{centering}

\caption{As waves propagates across space, the wavelength resulting from the
motion of the source will either stretch or contract depending on
its direction of motion. (Image was taken from http://en.wikipedia.org/wiki/File:Doppler\_effect\_diagrammatic.svg
on May 4, 2014.) \label{fig: Doppler Shift}}
\end{figure}


Additionally, the velocity of the observer and the source are determined
by the medium through which the waves are transmitted. The total Doppler
effect therefore results from the motion of the source, the observer,
and the medium and each of these effects must be analyzed separately.
To find the Doppler velocity of light waves, we can use Equation \ref{eq:doppler}
shown below. In this equation, $\Delta\nu=\nu_{obs}-\nu_{0}$ and
$c$ is the speed of light. However, for waves which do not encounter
much of a medium like light from point sources, only the relative
difference in velocities between the observer and the source have
to be considered. 

\begin{equation}
v_{doppler}=-\left(\frac{\Delta\nu}{\nu_{0}}\right)c\label{eq:doppler}
\end{equation}



\subsection{Galactic Coordinate Transformation}

The galactic coordinate system is an astronomical coordinate system
based on the position of the sun about the center of the Milky Way.
Galactic longitude ($l$) measures the angular distance of an object
eastward along the galactic equator from the galactic center. Galactic
latitude ($b$) measures the angular distance of an object above the
galactic equator with positive to the north and negative to the south.
These coordinates are an analogous to terrestrial longitude and latitude
and are usually measured in degrees.

\begin{figure}[H]
\begin{centering}
\includegraphics[scale=0.48]{/Users/ppkantorski/Documents/Radio_Astronomy/Lab_4/Pictures/Galactic_coordinates}
\par\end{centering}

\caption{In the galactic coordinate system, because of the large distances,
the Sun is placed at the origin. Galactic longitude ($l$) is measured
from the direction from the Sun to the center of the galaxy on the
galactic plane. Galactic latitude ($b$) measures the angle above
the galactic plane. Additionally the region between the sun's orbit
and the galactic center is known as the solar circle. (Image was taken
from http://en.wikipedia.org/wiki/File:Galactic\_coordinates.JPG on
May 5, 2014.)\label{fig: Galactic coordinates}}
\end{figure}


In spherical coordinates, the position of stars are often expressed
in right ascension and declination without use of a coordinate for
distance. Because celestial objects are extremely far way, astronomers
often have little to no information about their exact distances and
therefore use only the object's direction. Additionally, the direction
of fairly distant objects is the same for all observers and it is
convenient to specify this direction with the same coordinates for
every object. Using Equation \ref{eq:spherical} shown below, we can
convert from terrestrial latitude and longitude to spherical coordinates
($\alpha,\delta$).

\begin{equation}
\vec{x}=\left[\begin{array}{c}
\cos(lat)\times\cos(long)\\
\cos(lat)\times\sin(long)\\
\sin(lat)
\end{array}\right]\label{eq:spherical}
\end{equation}


The horizontal coordinate system is a celestial coordinate system
that utilizes an observer's local horizon as it's fundamental plane.
This coordinate system divides the sky into an upper hemisphere and
a lower hemisphere. In the upper hemisphere, objects and stars are
visible or within sight while in the lower hemisphere, the our view
of the sky is blocked by the Earth. The ``great circle'' separating
hemispheres is called celestial horizon or rational horizon and the
pole of the upper hemisphere is called the zenith. Using the spherical
coordinates ($\alpha,\delta$) and multiplying it by the rotational
matrix shown in Equation \ref{eq:sphere convert}, coordinates in
terrestrial latitude and longitude can be converted into celestial
coordinates ($az,alt$). Additionally, this depends on the observer's
terrestrial latitude $\phi$.

\begin{equation}
R_{(\alpha,\delta)\rightarrow(az,alt)}=\left[\begin{array}{ccc}
-\cos(LST)\times\sin(\phi) & -\sin(LST) & \cos(LST)\times\cos(\phi)\\
\sin(LST)\times\sin(\phi) & -\cos(LST) & -\sin(LST)\times\cos(\phi)\\
\cos(\phi) & 0 & \sin(\phi)
\end{array}\right]\label{eq:sphere convert}
\end{equation}


Because the Milky Way galaxy is extremely large and takes a long time
to complete a full rotation, galactic coordinates do not change much
on ``Earth-based'' scales of time. To convert terrestrial coordinates
to galactic, we can utilize the rotational matrix shown in Equation
\ref{eq:HA, Dec to l, b}. By multiplying spherical right ascension
and declination coordinates by the rotational matrix shown below in
Equation \ref{eq:HA, Dec to l, b}, we can convert ($\alpha,\delta$)
coordinates to ($l,b)$ coordinates. Since hour angle and declination
are ``Earth-based'' coordinates, this conversion will only depend
on the observer's terrestrial latitude.

\begin{equation}
R_{(\alpha,\delta)\rightarrow(l,b)}=\left[\begin{array}{ccc}
-0.054876 & -0.873437 & -0.483835\\
0.494109 & -0.444830 & 0.746982\\
-0.867666 & -0.198076 & 0.455984
\end{array}\right]\label{eq:HA, Dec to l, b}
\end{equation}



\section{Methods and Observations}


\subsection{Operating The Telescope}

To operate and control the 4.5m radio telescope at the Leuschner Observatory,
my partners and I programmed and remotely ran a Python script on Leuschner's
local computer via SSH. Since the H1-emission line is centered at
about 1420.4MHz and is present within a 12MHz Nyquist zone, we chose
to set our L.O. frequency to a range 12 Nyquist zones, or 144MHz,
away. We also added an additional 6MHz, half the nyquist frequency,
to center the H1-line. This made the center emission at a frequency
of about 1270.4MHz. For our observations, we also programmed the telescope
to take four separate measurements per pointing for the purpose of
calibrating our data. The four measurements taken include a sample
with the L.O.'s frequency set to 2MHz above the center frequency at
1272.4MHz, $10dB_{m}$, another with the frequency shifted 2MHz below
at 1268.4MHz, $10dB_{m}$, and two measurements using a noise diode
to add $T_{noise}=100K$ to the total system temperature $T_{sys}$
of each signal. When operating the telescope, we also had to ``home''
the telescope before each measurement to make sure that it was functioning
properly. Homing the dish resets the positioning drives of the telescopes
and returns the telescope to its default homing position (alt=90\textdegree{},
az=180\textdegree{}).


\subsection{Mapping Area}

The mapping area we covered for our observations was limited by the
amount of time give to our group to operate the telescope and the
pointing limits of the dish itself (see Figure \ref{fig:pointing limts}).

\begin{figure}[H]
\begin{centering}
\includegraphics[scale=0.46]{/Users/ppkantorski/Documents/Radio_Astronomy/Lab_4/Pictures/pointing_limits}
\par\end{centering}

\caption{This graph shows the pointing limits of the 4.5m radio telescope at
the Leuschner Observatory. Pointing limits result from the boundaries
of the telescope and it's location. \label{fig:pointing limts}}
\end{figure}


After calibrating data taken from May 2, 2014 and May 3, 2014, we
realizing that our coordinates were being processed incorrectly and
that the telescope did not always take valid measurements per pointing.
In the end, the total time range of our accurate observations was
diminished to May 4, 2014 from 17:48:37 - 20:41:31 UT. To compensate
for this, we set our galactic coordinates to $l=75^{\circ}$ and ranged
$b$ from $+20^{\circ}$ to $-20^{\circ}$, taking pointings at a
spacing of $2^{\circ}$. To ensure that we had all the data from each
pointing, we first set our coordinates to cover $4^{\circ}$ increments
within the range of $b$. After, we set them to cover $2^{\circ}$increments
and then another set of $4^{\circ}$increments, but skewed by $2^{\circ}$.
This ensured that we had at least one pointing for every position
in our $l$ and $b$ range. Additionally, we set the integration time
of each measurement to about 4 minutes per pointing or 8192 data points
per sample to ensure that our observation fell within our 3 hour observation
window.


\subsection{Signal Processing in Python}


\subsubsection{Manipulating Collected Data}

Observational data files collected from the radio telescope at The
Leuschner Observatory are stored in a '.log' format. To open these
files, I used the ``ReadSpec'' module provided on the UC Berkeley
Undergraduate Radio Lab GitHub page and then performed a mean calculation
on the data array. A mean was performed on the total dataset to reduce
the size and preserve the signal information but in a more manageable
format. To store data after calibration, I loaded the data into a
'.npz' file with the information on the entry numbers, dates of observation,
time stamps for each measurement, $b$ and $l$ coordinates, frequencies,
and temperatures. After calculating velocities and distances, I created
a separate program to create another '.npz' file to store additional
information on the Julian dates, right ascension and declination coordinates,
local standards of rest, velocities, and distances.


\subsubsection{Median Boxcar Smoothing}

Another signal processing method that I used to clean and filter out
bad data points from the raw data files is called ``Median Boxcar
Smoothing''. This is done by running a median across data within
a window of points before and after every data point. By taking the
median value within a window across all of the data, a new array of
``smooth'' data points positioned along the shape of the raw signal
will be produced. Because the number of elements in our data is approximately
8192, I made the radius of my boxcar 11 points.

In order to remove bad data points, I first subtracted out the smoothed
data from the raw data to normalize the signal about the frequency
axis. Then, by placing a cutoff range along the power axis above and
below four times the standard deviation of the total normalized signal,
I was able to find the position of each invalid point that deviates
significantly from the standard deviation value. Since both arrays
were of the same dimension, the outlier positions should be the same
for both datasets. Finally, I replaced the bad data in the raw data
with data points taken from the smoothed data.


\subsubsection{Polynomial Curve Fitting}

While calibrating the spectral lines, we also must remove the baseline
from the final product in order to isolate the H1-emission signal.
To do this, I used a 5th order polynomial curve to fit to the underlining
baseline shape of the data. Additionally, this was done Python using
the ``PolyFit'' function in the Numpy module.


\section{Data Analysis}


\subsection{Temperature Calibration}

The initial data points taken are from the radio telescope are recorded
in units of power versus frequency. To calibrate the data, I converted
the spectral lines into units of temperature with a script written
in Python. Before doing this, my partners and I programmed the telescope
to record four separate measurements per pointing with equal integration
times. As I previously mentioned, the four measurements include a
sample with the local oscillator's frequency set to 1272.4MHz at $10dB_{m}$,
another with the L.O. frequency shifted 4MHz below at 1268.4MHz, $10dB_{m}$,
and two measurements using a noise diode to add $T_{noise}=100K$
to the total system temperature $T_{sys}$ of each signal. To refer
to each of these signals respectfully, the subscript ``ONLINE''
was set for the higher frequency measurement and ``OFFLINE'' for
the lower one along with ``CALON'' or ``CALOFF'' for noise on
or off. The CALON measurements, as expected, increased the gain of
the signals. Additionally, after examining all four raw data measurements,
it was noticed that the gain of the lower 1268.4MHz signal was slightly
larger than the gain of the 1272.4MHz signal (see Figure \vref{fig:raw data}).

\begin{figure}[H]
\begin{centering}
\includegraphics[scale=0.5]{/Users/ppkantorski/Documents/Radio_Astronomy/Lab_4/Pictures/raw_data}
\par\end{centering}

\caption{This graph is of the four raw data measurements taken from May 4,
2014 from 19:44:09.18 - 19:48:02.82 UT at $b=-10^{\circ},$ $l=75^{\circ}$.
ONLINE measurements at 1272.4MHz are shown in black while OFFLINE
measurements at 1268.4MHz are shown in blue. CALON for each ONLINE
and OFFLINE measurement is the data shown in the same color above
while CALOFF is shown below. \label{fig:raw data}}
\end{figure}



\subsubsection{Calibration Process}

The first step to calibrate for the spectral lines is to find the
system temperature of each of the ONLINE and OFFLINE signals. As shown
in Equation \vref{eq:System Temperature}, by multiplying the noise
temperature with the proportionality between the total sum of the
CALON measurements and the sum of the difference between the CALON
and CALOFF measurements, the system temperature of each measurement
can be approximated.

\begin{equation}
T_{sys}=\left[\frac{\sum P_{j}^{CALON}}{\sum\left(P_{j}^{CALON}-P_{j}^{CALOFF}\right)}\right]\times T_{noise}\label{eq:System Temperature}
\end{equation}


After, the next step in calibration is to isolate the shape of the
ONLINE and OFFLINE, CALOFF signal. To do this, one of the measurements
have to be scaled down to the same level as the other. This can be
done by multiplying the OFFLINE dataset by the gain factor shown in
Equation \vref{eq:Gain Conversion}. Figure \vref{fig:gain fix} shows
the result of applying this equation to Figure \vref{fig:raw data}.

\begin{equation}
G=\frac{\sum P_{j}^{ONLINE,CALOFF}}{\sum P_{j}^{OFFLINE,CALOFF}}\label{eq:Gain Conversion}
\end{equation}


\begin{figure}[H]
\begin{centering}
\includegraphics[scale=0.5]{/Users/ppkantorski/Documents/Radio_Astronomy/Lab_4/Pictures/gain_fix}
\par\end{centering}

\caption{This graph is of the gain conversion and cleaning of data measurements
taken from May 4, 2014 from 19:44:09.18 - 19:48:02.82 UT at $b=-10^{\circ},$
$l=75^{\circ}$. ONLINE measurements at 1272.4MHz are shown in black
while OFFLINE measurements at 1268.4MHz are shown in blue. Signal
cleaning using the method of median boxcar smoothing is shown on the
graph on the right. As seen on the figures, the gain conversion is
not always perfect. However, this is offset is corrected isolating
``shape'' of the H1-signal and scaling it with the system temperature.\label{fig:gain fix}}
\end{figure}


After setting the ONLINE and OFFLINE, CALOFF measurements to the same
level, the next step is to extract the underlining signal ``shape''.
First, the OFFLINE measurements have to also be shifted back 4MHz
to the left so that they aligns with the ONLINE measurements. Next,
the section of the data including the H1-emission must be isolated
on each set along with the corresponding sections that do not have
any H1-emission present. Following Equation \vref{eq:HI Temperature},
by using the ``shape'' of the H1-signal and scaling it with the
system temperature, the data in terms of H1 temperature can finally
be calculated.

\begin{equation}
T_{ant,HI}^{ONLINE}(\nu)=\left[\frac{P_{+HI}^{ONLINE}}{P_{none}^{OFFLINE}}\right]\times T_{sys}^{ONLINE},\quad T_{ant,HI}^{OFFLINE}(\nu)=\left[\frac{P_{+HI}^{OFFLINE}}{P_{none}^{ONLINE}}\right]\times T_{sys}^{OFFLINE}\label{eq:HI Temperature}
\end{equation}


\begin{figure}[H]
\begin{centering}
\includegraphics[scale=0.5]{/Users/ppkantorski/Documents/Radio_Astronomy/Lab_4/Pictures/calibrate_raw}
\par\end{centering}

\caption{The graph shown above was made from data measurements taken from May
4, 2014 from 19:44:09.18 - 19:48:02.82 UT at $b=-10^{\circ},$ $l=75^{\circ}$.
ONLINE measurements at 1272.4MHz are shown in black while OFFLINE
measurements at 1268.4MHz are shown in blue. After applying Equation
\vref{eq:HI Temperature} to the section of data ranging from 1418.94
- 1421.86MHz, the signals become scaled to units of Kelvin. Because
there units in K cannot fall below 0K, a polynomial had to be fit
to the data in order to normalize it.\label{fig:calibrate_raw}}
\end{figure}


To convert both ONLINE and OFFLINE calibrated temperature spikes into
a final spectra, the underlining shape or ``base-line'' resulting
from the artifacts left in the signal must also be removed from the
datasets in order to normalize and preserve the H1-emission spectra.
This can be done by removing the H1-emission lines from the calibrated
temperature datasets and then fitting a polynomial function to the
shape of the base-line. Then, by subtracting the polynomial fit from
the original ONLINE and OFFLINE temperature calibrated datasets, the
H1-emission line on both datasets can be properly normalized. The
final step is to then combine and average the two isolated H1-emission
datasets (see Equation \vref{eq:Final HI Temperature}).

\begin{equation}
T_{ant,HI}(\nu)=\frac{1}{2}\left(T_{ant,HI}^{ONLINE}(\nu)+T_{ant,HI}^{OFFLINE}(\nu)\right)\label{eq:Final HI Temperature}
\end{equation}


\begin{figure}[H]
\begin{centering}
\includegraphics[scale=0.5]{/Users/ppkantorski/Documents/Radio_Astronomy/Lab_4/Pictures/calibrate_final}
\par\end{centering}

\caption{The graph shown above was made from data measurements taken from May
4, 2014 from 19:44:09.18 - 19:48:02.82 UT at $b=-10^{\circ},$ $l=75^{\circ}$
calibrated from power to units of Kelvin. After applying Equation
\vref{eq:Final HI Temperature} to the the difference between the
polynomial fit and the two datasets in Figure \vref{fig:calibrate_raw},
we get the final calibrated spectra.\label{fig:calibrate final}}
\end{figure}



\subsubsection{Error Propagation}

When taking a long integration and combining spectra, assuming each
spectra has statistically the same error $\Delta T_{sys}$, then the
error in the sum can be found as shown below in Equation \vref{eq:Error in Sum}.

\begin{equation}
(error\, in\, sum)^{2}=\left[\sum_{0}^{N-1}(\Delta T_{sys,n})^{2}\right]=N(\Delta T_{sys})^{2}\label{eq:Error in Sum}
\end{equation}


\begin{equation}
(error\, in\, average)^{2}=\frac{\left(\Delta T_{sys}\right)^{2}}{N}\label{eq:Error in Avg}
\end{equation}


The average value of a dataset is equal to the sum divided by $N$.
This makes the error associated with averaging squared approximately
equal to the error in the sum squared divided by $N^{2}$. Thus the
error in averaging is shown by Equation \vref{eq:Error in Avg}. After
taking the square root of the equation, we can see how how the noise
in the data measurements decreases with the square root of $N$, the
number of times that a measurement is repeated. Applying this equation
to the general form of Equation \vref{eq:HI Temperature}, we get
the following expression.

\begin{equation}
\delta T_{sys}(\nu)^{2}=\left[\left(\frac{\delta P_{j}^{ONLINE,CALOFF}}{P_{j}^{ONLINE,CALOFF}}\right)^{2}+\left(\frac{\delta P_{j}^{OFFLINE,CALOFF}}{P_{j}^{OFFLINE,CALOFF}}\right)^{2}\right](T_{sys})^{2}\label{eq:Error of Seperate}
\end{equation}


Since the statistical properties of $P_{j}^{ONLINE,CALOFF}$and $P_{j}^{OFFLINE,CALOFF}$are
the same, we can simplify the expression and combine the contributions
of the two spectra (see Equation \vref{eq:Error Total}). Therefore,
the noise in $\delta T_{sys}(\nu)^{2}$ increases by a multiplicative
factor of 2 or the noise in $T_{sys}(\nu)$ by a factor $\sqrt{2}$.
Although the ONLINE measurements contain the desired signal information,
it is false to assume that spending more time collecting ONLINE than
the OFFLINE measurements will optimize the accuracy of the data calibration.
The best results are obtained by spending an equal amount of time
on each pointing. This is because equal amounts of integration time
produces statistically equal amounts of signal noise per pointing.
Furthermore, since there is no way obvious way around inducing some
signal noise into our data calibration, it is impossible to generate
a completely ``noise-free'' spectrum.

\begin{equation}
\delta T_{sys}(\nu)^{2}=2\left(\frac{\delta P_{j}^{ONLINE,CALOFF}}{P_{j}^{ONLINE,CALOFF}}\right)^{2}(T_{sys})^{2}\label{eq:Error Total}
\end{equation}



\subsection{Cloud Velocities}

Since the HI-line frequency is centered about 1420.4MHz, when the
line appears at different frequencies, we can conclude that the signals
must have been Doppler shifted. Using Equation \vref{eq:doppler},
the velocities corresponding to each frequency can be calculated.
However, the return accurate Doppler velocity values, we much also
correct our measurements to the Local Standard of Rest (LSR). The
LSR is a velocity that accounts for the movement of our solar system
about the galactic center. In our observation, our LSR correction
ranged from -28.9693km/s to -35.7493km/s. Furthermore, to find the
velocities corresponding to our line of sight, we must also divide
our velocity by the cosine of the inclination angle $b$ (See Equation
\vref{eq:correct velocity}).

\begin{equation}
v_{measured}=\left(v_{doppler}-LSR\right)\left(\frac{1}{\cos\left(b\right)}\right)\label{eq:correct velocity}
\end{equation}


Figure \vref{fig:velocity vs temp} shows the result of converting
the frequency axis of our dataset into velocities plotted against
temperature. For positions within the solar circle that our solar
system is approaching (see Figure \vref{fig: Galactic coordinates}),
velocities are red shifted while positions outside are blue shifted.
On the contrary, for positions within the solar circle that our solar
system is receding from, the velocities are blue shifted while the
positions outside are red shifted. This red and blue shift is the
result of our location direction of motion within the Milky Way. 

\begin{figure}[H]
\begin{centering}
\includegraphics[scale=0.5]{/Users/ppkantorski/Documents/Radio_Astronomy/Lab_4/Pictures/velocity_vs_temp}
\par\end{centering}

\caption{This graph is of all the calibrated data from the May 4, 2014 observation
in terms of velocity (km/s) and temperature (K). The data points on
the left side of the dashed line are signals that has been red shifted
while data points to the right are blue shifted. \label{fig:velocity vs temp}}
\end{figure}



\subsection{Distance Calculations}

To convert our data from velocities into the associated corresponding
distances, the angle $l$, velocity of the Sun about the galactic
center $v_{0}\approx220km/s$, and the distance between the two $D_{0}\approx8kpc$
can be utilized. If we form a right triangle between the Earth, the
Sun, and positions outside the solar circle, then the angle $\alpha$
can be thought of as the additional angle behind the Sun required
to complete $90^{\circ}$ (see Equation \vref{eq:angular conversion factor}).
Additionally, since we know the Sun's velocity, the corrected Doppler
velocities, and galactic longitude $l$, $\alpha$ can be calculated
for each dataset.

\begin{equation}
\alpha=\arcsin\left(\frac{v_{measured}}{v_{0}}-\sin\left(l\right)\right)+l\label{eq:angular conversion factor}
\end{equation}


Using angle $\alpha$, the distance between the Sun and the galactic
circle $D_{0}$, and some clever geometry, we can derive a distance
equation for our measurements (see Equation \vref{eq:distance eqn}). 

\begin{equation}
D=D_{0}\left[\sin\left(90^{\circ}-l\right)+\sin\left(l\right)\tan\left(90^{\circ}-\alpha+l\right)\right]\label{eq:distance eqn}
\end{equation}


After converting the velocities into units of distance, my partners
and I also solved for the heights above and below the galactic plane
by using the distance $D$ and galactic latitude angle $b$ (see Equation
\vref{eq:height eqn}).

\begin{equation}
H=D\times\tan\left(b\right)\label{eq:height eqn}
\end{equation}


\begin{figure}[H]
\begin{centering}
\includegraphics[scale=0.56]{/Users/ppkantorski/Documents/Radio_Astronomy/Lab_4/Pictures/dist_height_temp}
\par\end{centering}

\caption{This graph is of all the calibrated data from the May 4, 2014 observation
in terms of distance (kpc) and temperature (K) on the left and distance
(kpc) and height (kpc) above the galactic plane on the right. As the
distances between us and what we are observing increases, we begin
to see a drop in the overall temperature. Values ranging from 0 -
3kpc are emissions that has been red shifted. Because of the limited
amount of time, my partners and I decided to neglect these values.
Additionally, because measurements were taken spaced out in $b$ by
$2^{\circ}$, the resolution of our dataset was fairly limited. \label{fig:dist vs velocity & dist vs height}}
\end{figure}



\subsection{Imaging the Galactic Warp}


\subsubsection{Generating the Image}

Using the final calibrated data, we were able to construct an image.
By averaging degenerate datasets, convolving them with a gaussian
filter, and zeroing values with intensities below 0K, we were able
to map intensity to the second plot shown on Figure \vref{fig:dist vs velocity & dist vs height}.
To do this, we first constructed a new image array with x values ranging
from 0kpc to 25kpc and y values from -10kpc to 10kpc. Treating each
square of information on this picture as a single pixel, the final
picture came out to a resolution of roughly 0.1kpc/pix. After, data
points were then divided by the resolution to fill every pixel on
the array with information with intensities mapped to each pixel.

\begin{equation}
kernal=\exp-\left[\frac{x^{2}}{2\sigma_{x}^{2}}+\frac{y^{2}}{2\sigma_{y}^{2}}\right]\label{eq:kernal}
\end{equation}


The next step was to define weights for the image in order to keep
track of the number of data points per pixel and create Gaussian kernel
function to convolve or smooth the intensities (See Equation \vref{eq:kernal}).
By convolving our image and then dividing it by the weights, we were
able to essentially average our data points, bringing them to the
corrected levels of intensities. To ensure that the image only convolves
meaningful data points, if the pixel of information was not within
a range of 3 pixels, the weights were set to 0K.

\begin{figure}[H]
\begin{centering}
\includegraphics[scale=0.42]{/Users/ppkantorski/Documents/Radio_Astronomy/Lab_4/Pictures/final_image_fix}
\par\end{centering}

\caption{The graph above is a plot all of the data collected from the May 4,
2014, 17:48:37 - 20:41:31 UT observation. Information is plotted in
terms distance from our solar system, height above the galactic plane,
and temperature intensity in units of Kelvin. on the right. As the
distances between us and what we are observing increases, we begin
to see a drop in the overall temperature. Additionally, because we
took measurements spaced out in $b$ by $2^{\circ}$, the resolution
of our dataset shown on the right was limited. For prosthetic purposes,
red shifted data points or points between 0-3kpc were filled in with
the intensity along the border between red and blue shifts. \label{fig:final image}}
\end{figure}



\subsubsection{Basic Photometric Implications}

The optical depth of the H1-emissions along the galactic plane can
be taken from the application of Equation \vref{eq:optical depth},
where $T_{actual}$ is temperature intensity, $\tau$ is the optical
thickness, $\sigma_{area}$ is the cross-sectional area of the H1
emissions, and $N_{emission}$ is the column density or amount of
particles per unit of volume.

\begin{equation}
(optical\, depth)=T_{actual}e^{-\tau},\quad\tau=\sigma_{area}N_{emission}\label{eq:optical depth}
\end{equation}


Small values of optical thickness correspond to more information within
the visible line of sight. In the case when the optical thickness
$\tau\ll1$, the optical depth reduces to the form $T_{actual}\sigma_{area}N_{emission}$.
So, for optically thin regions, the optical depth is larger. When
viewing Figure \vref{fig:final image}, the regions through which
H1 temperatures are plotted are optically thin while the lower temperature
regions are optically thicker or experience more extinction as a result
of dense H1 gas. The region around a distance of 10kpc looks like
it may be extinction resulting from the spiral arms of our galaxy.
Furthermore, because of the uneven distribution of H1 gas and uneven
distributions above and below the galactic plane, Figure \vref{fig:final image}
supports the concept of galactic warping in the Milky Way. However,
additional data collections may be required before a solid conclusion
can be derived.


\section{Conclusion}


\subsection{Troubleshooting}

Several problems were encountered while trying to carry out this lab.
The main issue was with data collection and resulted from an incorrect
application of coordinate transformations. This was caught by me when
I was calibrating data sets taken from May 1, 2014 to May 3, 2014.
If the telescope reaches the boundary of its pointing limits, the
H1-line ends up skewed to positions where it should not be. After
viewing this happening on a coordinate that was known to be within
the pointing limits of the telescope, my partners and I immediately
knew something was wrong. After fixing the rotational matrices in
our observation code, my partners and I were limited to a 3 hour window
on May 4, 2014 for data collection. Other groups performing experiments
using the 4.5m radio telescope at the Leuschner Observatory were given
more time for their observations based on necessity to complete their
report. Since our experiment did not require as much data, the window
of observation given to us was already small.

Another issue we noticed was that some of our observations were returning
empty values. To ensure this problem did not happen again, we took
multiple vertical profiles for our last observation, reducing the
resolution of our final image. Nevertheless, with 3 hours of observation
time, we were still able to show evidence to support the galactic
warp of the Milky Way galaxy.


\subsection{Results}

With more vertical profiles of the Milky Way at a higher resolution,
my partners and I would have been able to make a more detailed image
of the galactic warp of the Milky Way galaxy. By plotting multiple
profiles and mapping it to the vertical displacement of the H1-emissions,
a picture like Figure \vref{fig:galactic warp} could have been constructed.
Regardless, I personally am fairly satisfied with the image we constructed
in the limited amount of observation time we were given. 

\begin{figure}[h]
\begin{centering}
\includegraphics[scale=0.52]{/Users/ppkantorski/Documents/Radio_Astronomy/Lab_4/Pictures/galactic_warp}
\par\end{centering}

\caption{Final galactic warp image created by the undergraduate radio astronomy
group Gamma (Andrew, Suneet, James, Nhung) made from data collected
using the 4.5m radio telescope at the Leuschner Observatory in 2006.
(Image was taken from http://ugastro.berkeley.edu/radio/leuschner/previous\_results/gplane\_warp.pdf
on May 6, 2014.). \label{fig:galactic warp}}
\end{figure}


Radio telescopes are extremely powerful instruments capable of detecting
emissions beyond what is observable from the optical wavelengths.
By taking advantage of the 21cm line of H1-gas, regions of emission
can be mapped to temperature intensities, doppler velocities, and
distances. Ultimately, this lab has taught how to properly calibrate
H1-emission lines, operate the radio telescope at the Leuschner Observatory,
derive coordinates for detailed observations, and map plots to positions
along regions of space to H1 temperature intensities. Since I want
to enter the field of astronomy in the future, the undergraduate radio
astronomy course offered at the University of California, Berkeley
has been one of the most interesting and important classes for me
and my desire to pursue research. 


\section{Acknowledgements and References}
\begin{enumerate}
\item Lab Partners: Sahana Kumar, Daniel Krishan, and Marta Pastore.
\item Additional help was also proved by Baylee Bordwell and Eduardo Herrera.
\item Undergraduate radio lab. (2014, May 5). Retrieved May 5, 2014 from


https://casper.berkeley.edu/astrobaki/index.php/Undergraduate\_Radio\_Lab

\item Hare, J. (2005, September 23). Leuschner Observatory. Retrieved May
5, 2014, from


http://ugastro.berkeley.edu/leuschner/old\_site/index.html\end{enumerate}

\end{document}
